\chapter{Field Extensions}

\begin{center}\itshape
With a ruler and a compass you can bisect any angle.\\
Can you trisect an angle of $60^\circ$? Can you build a cube of twice the volume of a given one?\\
For two thousand years nobody could, and nobody could prove that it was impossible.
\end{center}

You can bisect a $60^\circ$ angle with straightedge and compass. Can you cut it into three equal $20^\circ$ angles with the same tools?

Let $u=2\cos20^\circ$. The triple-angle identity gives
\[
u^3-3u-1=0.
\]
The polynomial $X^3-3X-1$ has no rational root: the only possible integer roots are $1$ and $-1$, and neither works. A cubic without a rational root is irreducible over $\Q$. So $u$ has degree $3$ over $\Q$. The relevance of that number will become clear once we understand which degrees a straightedge-and-compass construction can produce.

A field extension $L/K$ means that the field $L$ contains the smaller field $K$. We can then view $L$ as a vector space over $K$, and write its dimension as $[L:K]$. If $K\subseteq M\subseteq L$, those dimensions multiply:
\[
[L:K]=[L:M][M:K].
\]
This is the tower law. Since intersecting lines and circles only requires solving linear and quadratic equations, constructible numbers fit into towers whose successive degrees are at most $2$. Their degrees over $\Q$ must therefore be powers of $2$, ruling out $u$.

The same machinery prepares the next questions. What field must we add so that a polynomial splits completely into linear factors? When is that field determined up to isomorphism? These are the fields whose symmetries Galois theory will study.

\section{Degree and the tower law}

\begin{definition}
A \emph{field extension} $L/K$ is a field $L$ together with a subfield $K$. Then $L$ is a vector space over $K$ (add in $L$, multiply by elements of $K$); its dimension is the \emph{degree} $[L:K]\in\{1,2,\dots\}\cup\{\infty\}$. The extension is \emph{finite} if $[L:K]<\infty$. A field $M$ with $K\subseteq M\subseteq L$ is an \emph{intermediate field}.
\end{definition}

For instance $[\C:\R]=2$, with basis $1,i$. The degree $[\R:\Q]$ is infinite: a $\Q$-vector space of finite dimension $n$ is in bijection with $\Q^n$, which is countable, while $\R$ is not.

\begin{theorem}[tower law]\label{thm:tower}
Let $K\subseteq L\subseteq M$ be fields. Then $M/K$ is finite iff $M/L$ and $L/K$ are, and then
\[
[M:K]=[M:L]\,[L:K].
\]
More precisely, if $(e_i)_{i\in I}$ is a basis of $L$ over $K$ and $(f_j)_{j\in J}$ a basis of $M$ over $L$, then $(e_if_j)_{(i,j)\in I\times J}$ is a basis of $M$ over $K$.
\end{theorem}
\begin{proof}
We first show that the products $e_if_j$ span $M$ over $K$. Take $z\in M$. Since $(f_j)$ is a basis over $L$, we may write
\[
z=\sum_j\lambda_jf_j,\qquad\lambda_j\in L,
\]
with only finitely many non-zero coefficients. Each $\lambda_j$ has its own expansion in the $K$-basis $(e_i)$:
\[
\lambda_j=\sum_i\mu_{ij}e_i,\qquad\mu_{ij}\in K.
\]
Substitution gives $z=\sum_{i,j}\mu_{ij}e_if_j$. So the products span.

Next suppose that a finite $K$-linear combination of these products is zero:
\[
0=\sum_{i,j}c_{ij}e_if_j
=\sum_j\Bigl(\sum_i c_{ij}e_i\Bigr)f_j.
\]
Each parenthesis belongs to $L$. The $f_j$ are linearly independent over $L$, so every parenthesis must vanish. Now use the independence of the $e_i$ over $K$: each $c_{ij}$ must be zero. The products are therefore linearly independent.

They form a basis of $M$ over $K$, indexed by $I\times J$. Both index sets are non-empty because a field has dimension at least one over any subfield. Consequently $I\times J$ is finite exactly when both $I$ and $J$ are finite. In that case,
\[
[M:K]=|I\times J|=|I|\,|J|=[L:K]\,[M:L].
\]
\end{proof}

\paragraph{Why this works.}
A basis lets us count coordinates. If $L$ has basis $e_1,\ldots,e_r$ over $K$, and $M$ has basis $f_1,\ldots,f_s$ over $L$, then a coordinate in $M$ first chooses one of the $s$ vectors $f_j$, and its coefficient needs $r$ coordinates from $K$. We therefore get $rs$ coordinates altogether. The independence part of the proof is essential: merely writing down $rs$ spanning vectors does not show that the extension has degree $rs$. For example, $\Q(\sqrt2,\sqrt3)$ has degree $4$ over $\Q$, using the basis $1,\sqrt2,\sqrt3,\sqrt6$, because $\Q(\sqrt2,\sqrt3)$ has degree $2$ over $\Q(\sqrt2)$ and $\Q(\sqrt2)$ has degree $2$ over $\Q$.

\begin{corollary}\label{cor:towerdiv}
If $K\subseteq L\subseteq M$ and $M/K$ is finite, then $[L:K]$ and $[M:L]$ divide $[M:K]$. In particular, if $[M:K]$ is prime, there is no intermediate field other than $K$ and $M$.
\end{corollary}
\begin{proof}
The first sentence is the tower law. If $[M:K]=p$ is prime, $[L:K]\in\{1,p\}$; if it is $1$ then $L=K$, and if it is $p$ then $[M:L]=1$ and $L=M$.
\end{proof}

Degrees measure the size of an extension, but a concrete extension often begins with one number: $\sqrt2$, $\sqrt[3]2$, or a complex root of a polynomial. The polynomial of smallest degree that vanishes at that number tells us its degree over the base field.

\section{Algebraic elements}

Let $L/K$ be an extension and $\alpha\in L$. We write $K[\alpha]=\{P(\alpha):P\in K[X]\}$, the smallest subring of $L$ containing $K$ and $\alpha$, and $K(\alpha)$ for the smallest subfield containing them. More generally $K(\alpha_1,\dots,\alpha_n)$ is the smallest subfield of $L$ containing $K$ and the $\alpha_i$; it equals $K(\alpha_1,\dots,\alpha_{n-1})(\alpha_n)$, since both are the smallest subfield containing $K$ and all the $\alpha_i$.

\begin{definition}
$\alpha$ is \emph{algebraic} over $K$ if $P(\alpha)=0$ for some non-zero $P\in K[X]$, and \emph{transcendental} otherwise. An extension is \emph{algebraic} if all its elements are.
\end{definition}

\begin{proposition}[minimal polynomial]\label{prop:minpoly}
Let $\alpha\in L$ be algebraic over $K$. There is a unique monic polynomial $m_\alpha\in K[X]$ such that, for $P\in K[X]$,
\[
P(\alpha)=0\iff m_\alpha\mid P.
\]
It is irreducible, and it is the only monic irreducible polynomial of $K[X]$ with root $\alpha$. Its degree is the \emph{degree of $\alpha$ over $K$}.
\end{proposition}
\begin{proof}
The evaluation map $\operatorname{ev}_\alpha:K[X]\to L$, $P\mapsto P(\alpha)$, is a ring homomorphism, so its kernel $I$ is an ideal, non-zero as $\alpha$ is algebraic. $K[X]$ is a PID (Chapter~6), so $I=(m)$, and $m$ can be taken monic; two monic generators divide each other, so they are equal. This is the displayed property.

If $m=fg$ with $f,g\in K[X]$, then $f(\alpha)g(\alpha)=0$ in the field $L$, so say $f(\alpha)=0$. Then $m\mid f$, so $\deg f\ge\deg m$ and $g$ is a constant: $m$ is irreducible. Finally, if $P$ is monic irreducible with $P(\alpha)=0$, then $m\mid P$, and as $P$ is irreducible and $m$ is not constant, $P$ is $m$ times a constant, which is $1$ because both are monic.
\end{proof}

\begin{theorem}[simple extensions]\label{thm:simple}
Let $\alpha\in L$ be algebraic over $K$, of degree $d$. Then
\[
K[X]/(m_\alpha)\;\xrightarrow{\ \sim\ }\;K[\alpha],\qquad P+(m_\alpha)\mapsto P(\alpha),
\]
is a ring isomorphism, $K[\alpha]=K(\alpha)$ is a field, and $1,\alpha,\dots,\alpha^{d-1}$ is a basis of $K(\alpha)$ over $K$. In particular $[K(\alpha):K]=d$. If $\alpha$ is transcendental, $[K(\alpha):K]=\infty$.
\end{theorem}
\begin{proof}
The image of $\operatorname{ev}_\alpha$ is $K[\alpha]$ and its kernel is $(m_\alpha)$, so the first isomorphism theorem (Proposition~\ref{prop:quotring}) gives the isomorphism. In the PID $K[X]$, the ideal generated by an irreducible element is maximal (proof of Theorem~\ref{thm:edpidufd}), so $K[X]/(m_\alpha)$ is a field (Proposition~\ref{prop:primemax}), and so is $K[\alpha]$. A subfield of $L$ containing $K$ and $\alpha$ contains $K[\alpha]$, so $K[\alpha]=K(\alpha)$.

\emph{Basis.} Every $P\in K[X]$ is $Qm_\alpha+R$ with $\deg R<d$ (Euclidean division), so $P(\alpha)=R(\alpha)$ is a $K$-combination of $1,\alpha,\dots,\alpha^{d-1}$. If $\sum_{k<d}c_k\alpha^k=0$, the polynomial $\sum_{k<d}c_kX^k$ is a multiple of $m_\alpha$ of degree $<d$, hence $0$.

If $\alpha$ is transcendental, the powers $1,\alpha,\alpha^2,\dots$ are independent over $K$ (a relation would be a polynomial vanishing at $\alpha$), so $K(\alpha)$ has infinite dimension.
\end{proof}

\begin{corollary}\label{cor:finitealg}
A finite extension $L/K$ is algebraic, and the degree of each $\alpha\in L$ divides $[L:K]$. Conversely, if $\alpha_1,\dots,\alpha_n$ are algebraic over $K$, then $K(\alpha_1,\dots,\alpha_n)/K$ is finite.
\end{corollary}
\begin{proof}
If $[L:K]=n$, the $n+1$ elements $1,\alpha,\dots,\alpha^n$ are linearly dependent over $K$, which is a non-zero polynomial of degree $\le n$ vanishing at $\alpha$. Then $K\subseteq K(\alpha)\subseteq L$ and the tower law give $[K(\alpha):K]\mid[L:K]$.

For the converse, $\alpha_k$ is algebraic over $K$, hence over $K_{k-1}=K(\alpha_1,\dots,\alpha_{k-1})$ (the same polynomial works), so $[K_k:K_{k-1}]$ is finite by Theorem~\ref{thm:simple}; multiply along the tower.
\end{proof}

Theorem~\ref{thm:simple} also says how to compute in $K(\alpha)$: an element is a polynomial in $\alpha$ of degree $<d$, products are reduced modulo $m_\alpha$, and inverses come from Bézout. In $\Q(\sqrt2)$, for instance, $m=X^2-2$ and $1/(1+\sqrt2)=\sqrt2-1$ because $(X+1)(X-1)=X^2-1\equiv1\pmod{X^2-2}$. This is how computer algebra systems handle algebraic numbers exactly.

\begin{think}
What is the degree of $\sqrt2+\sqrt3$ over $\Q$? It is a sum of two numbers of degree $2$; is it of degree $2$, or more?
\end{think}

\begin{example}\label{ex:sqrt2sqrt3}
Let $\alpha=\sqrt2+\sqrt3$. Then $\alpha^2=5+2\sqrt6$, so $(\alpha^2-5)^2=24$ and $\alpha$ is a root of $X^4-10X^2+1$. Its degree is $4$:
\begin{enumerate}[label=(\arabic*)]
\item $\sqrt3\notin\Q(\sqrt2)$. Otherwise $\sqrt3=a+b\sqrt2$ with $a,b\in\Q$, and squaring gives $3=a^2+2b^2+2ab\sqrt2$, so $ab=0$ (as $\sqrt2\notin\Q$). If $b=0$, $\sqrt3=a\in\Q$; if $a=0$, $3=2b^2$, so $\sqrt{3/2}=|b|\in\Q$. Both are false. Hence $X^2-3$ has no root in $\Q(\sqrt2)$, it is irreducible there, and by the tower law $[\Q(\sqrt2,\sqrt3):\Q]=2\cdot2=4$.
\item $\Q(\alpha)=\Q(\sqrt2,\sqrt3)$. Expanding, $\alpha^3=11\sqrt2+9\sqrt3$, so
\[
\sqrt2=\tfrac12(\alpha^3-9\alpha),\qquad\sqrt3=\alpha-\sqrt2
\]
lie in $\Q(\alpha)$; the reverse inclusion is clear.
\end{enumerate}
So $[\Q(\alpha):\Q]=4$, and $X^4-10X^2+1$, monic of degree $4$ with root $\alpha$, is $m_\alpha$. Figure~\ref{fig:pentagon} (right) shows the three intermediate fields of degree $2$; Chapter~10 proves that there are no others.
\end{example}

The example is typical: a field generated by several algebraic numbers is often generated by a single one, a \emph{primitive element}. Over $\Q$ this is always true (Theorem~\ref{thm:primitiveelement}).

Adjoining one root may leave several others missing. For the symmetries of a polynomial to act on all its roots, we need a field containing all of them. Such a splitting field is the stage on which Galois theory will work.

\section{Splitting fields}\label{sec:splitting}

So far the big field $L$ was given. Often only $K$ and a polynomial are given, and we want a field in which the polynomial has roots. Kronecker's idea is to \emph{define} a root as the class of $X$, as $\C$ was defined as $\R[X]/(X^2+1)$ in Chapter~6.

\begin{theorem}[Kronecker]\label{thm:kronecker}
Let $P\in K[X]$ be irreducible of degree $d$. Then $L=K[X]/(P)$ is a field; $c\mapsto c+(P)$ identifies $K$ with a subfield of $L$; the class $x=X+(P)$ is a root of $P$ in $L$; and $L=K(x)$, $[L:K]=d$. Consequently every non-constant polynomial over $K$ has a root in some finite extension of $K$.
\end{theorem}
\begin{proof}
$(P)$ is maximal (as in Theorem~\ref{thm:simple}), so $L$ is a field. The map $c\mapsto c+(P)$ is a ring homomorphism from a field, so it is injective (its kernel is an ideal of $K$ not containing $1$, hence $0$). If $P=\sum a_kX^k$, then $P(x)=\sum a_kx^k$ is the class of $\sum a_kX^k=P$, which is $0$ in $L$. Every element of $L$ is the class of a polynomial, i.e.\ a polynomial in $x$, so $L=K[x]=K(x)$; $P$ is monic irreducible up to a constant with root $x$, so $[L:K]=d$ by Proposition~\ref{prop:minpoly} and Theorem~\ref{thm:simple}. For the last sentence, apply this to an irreducible factor of the given polynomial.
\end{proof}

\begin{definition}
Let $P\in K[X]$ have degree $n\ge1$. A \emph{splitting field} of $P$ over $K$ is an extension $L/K$ such that $P=c(X-\alpha_1)\cdots(X-\alpha_n)$ in $L[X]$ and $L=K(\alpha_1,\dots,\alpha_n)$.
\end{definition}

\begin{proposition}\label{prop:splitexists}
Every $P\in K[X]$ of degree $n\ge1$ has a splitting field $L$, with $[L:K]\le n!$.
\end{proposition}
\begin{proof}
Induction on $n$, for all fields $K$ at once. If $n=1$, $L=K$. Otherwise, Theorem~\ref{thm:kronecker} gives $K_1=K(\alpha_1)$ with $P(\alpha_1)=0$ and $[K_1:K]\le n$. In $K_1[X]$, $P=(X-\alpha_1)Q$ with $\deg Q=n-1$. By induction $Q$ has a splitting field $L=K_1(\alpha_2,\dots,\alpha_n)$ over $K_1$ with $[L:K_1]\le(n-1)!$. Then $P$ splits in $L$, $L=K(\alpha_1,\dots,\alpha_n)$, and $[L:K]\le n\cdot(n-1)!$.
\end{proof}

To compare two splitting fields we need to transport fields along isomorphisms. If $\sigma:K\to K'$ is a field isomorphism and $P=\sum a_kX^k\in K[X]$, write $\sigma P=\sum\sigma(a_k)X^k$. Then $P\mapsto\sigma P$ is a ring isomorphism $K[X]\to K'[X]$, so it preserves divisibility and irreducibility.

\begin{lemma}\label{lem:extend1}
Let $\sigma:K\to K'$ be an isomorphism, $Q\in K[X]$ irreducible, $\alpha$ a root of $Q$ in an extension of $K$ and $\beta$ a root of $\sigma Q$ in an extension of $K'$. There is exactly one isomorphism $\tau:K(\alpha)\to K'(\beta)$ that extends $\sigma$ and sends $\alpha$ to $\beta$.
\end{lemma}
\begin{proof}
Up to a constant, $Q$ is the minimal polynomial of $\alpha$ and $\sigma Q$ that of $\beta$ (Proposition~\ref{prop:minpoly}). The isomorphism $P\mapsto\sigma P$ maps the ideal $(Q)$ onto $(\sigma Q)$, so it induces an isomorphism $K[X]/(Q)\to K'[X]/(\sigma Q)$. Compose with the isomorphisms of Theorem~\ref{thm:simple}:
\[
K(\alpha)\cong K[X]/(Q)\cong K'[X]/(\sigma Q)\cong K'(\beta),
\]
which sends $P(\alpha)$ to $(\sigma P)(\beta)$; in particular it extends $\sigma$ and sends $\alpha$ to $\beta$. Uniqueness: an extension of $\sigma$ with $\alpha\mapsto\beta$ must send $\sum a_k\alpha^k$ to $\sum\sigma(a_k)\beta^k$, and every element of $K(\alpha)$ has this form.
\end{proof}

In characteristic $0$, irreducible polynomials have no repeated roots. This innocent fact is what makes the counts of Galois theory exact.

\begin{lemma}\label{lem:sep}
Let $K$ have characteristic $0$ and $Q\in K[X]$ be irreducible. Then $Q$ has distinct roots in every extension of $K$.
\end{lemma}
\begin{proof}
Let $\deg Q=d\ge1$ and $Q=\sum_{k\le d}a_kX^k$. The derivative $Q'=\sum ka_kX^{k-1}$ has leading coefficient $da_d\ne0$ (characteristic $0$), so $Q'\ne0$ and $\deg Q'=d-1$. A gcd of $Q$ and $Q'$ divides the irreducible $Q$, so it is $1$ or $Q$ up to a constant; $Q$ does not divide $Q'$ by degree, so the gcd is $1$ and $UQ+VQ'=1$ for some $U,V\in K[X]$ (Bézout in the PID $K[X]$). If $Q=(X-\alpha)^2S$ in $L[X]$, then $Q'=2(X-\alpha)S+(X-\alpha)^2S'$, so $Q(\alpha)=Q'(\alpha)=0$, and evaluating $UQ+VQ'=1$ at $\alpha$ gives $0=1$.
\end{proof}

\begin{remark}
In characteristic $p$ this can fail: over the field $\mathbb F_p(t)$ of rational fractions, $X^p-t$ is irreducible but equals $(X-\theta)^p$ in a splitting field, because $(a+b)^p=a^p+b^p$ in characteristic $p$. Over finite fields it does not happen (Chapter~9).
\end{remark}

\begin{theorem}[uniqueness of splitting fields]\label{thm:splitunique}
Let $\sigma:K\to K'$ be an isomorphism, $P\in K[X]$ of degree $\ge1$, $L$ a splitting field of $P$ over $K$ and $L'$ a splitting field of $\sigma P$ over $K'$. Then $\sigma$ extends to an isomorphism $L\to L'$, and the number of such extensions is at most $[L:K]$, with equality if $\sigma P$ has distinct roots in $L'$.
In particular a splitting field is unique up to isomorphism, and if $P$ has distinct roots in $L$, then $L$ has exactly $[L:K]$ automorphisms that fix $K$ pointwise.
\end{theorem}
\begin{proof}
Induction on $[L:K]$, which is finite by Corollary~\ref{cor:finitealg}. If $[L:K]=1$, $P$ splits over $K$, so $\sigma P$ splits over $K'$ (apply $\sigma$ to the factorisation), $L'=K'$, and $\sigma$ itself is the only extension.

If $[L:K]>1$, some root $\alpha\in L$ of $P$ is not in $K$. Let $Q$ be its minimal polynomial, of degree $d>1$; $Q\mid P$, so $\sigma Q\mid\sigma P$ and $\sigma Q$ splits in $L'$, with $e\le d$ distinct roots $\beta_1,\dots,\beta_e$; if $\sigma P$ has distinct roots, so has its factor $\sigma Q$, and $e=d$.

Any extension $\tau:L\to L'$ of $\sigma$ satisfies $(\sigma Q)(\tau(\alpha))=\tau(Q(\alpha))=0$, so $\tau(\alpha)=\beta_j$ for some $j$, and then $\tau$ restricts on $K(\alpha)$ to the isomorphism $\sigma_j:K(\alpha)\to K'(\beta_j)$ of Lemma~\ref{lem:extend1}. Now $L$ is a splitting field of $P$ over $K(\alpha)$ (it is generated by the roots of $P$ over $K$, hence over $K(\alpha)$), $L'$ is one of $\sigma_jP=\sigma P$ over $K'(\beta_j)$, and $[L:K(\alpha)]=[L:K]/d<[L:K]$. By induction, each $\sigma_j$ has at least one and at most $[L:K]/d$ extensions $L\to L'$, exactly $[L:K]/d$ if $\sigma P$ has distinct roots. Each extension of $\sigma$ arises from exactly one $j$, so there are between $1$ and $e\cdot[L:K]/d\le[L:K]$ of them, and exactly $[L:K]$ when the roots are distinct.

For the last sentence take $K'=K$, $L'=L$ and $\sigma=\mathrm{id}$.
\end{proof}

In characteristic $0$ the hypothesis ``distinct roots'' holds as soon as $P$ is a product of distinct monic irreducibles: each has distinct roots (Lemma~\ref{lem:sep}), and two of them cannot share a root, since both would be its minimal polynomial. The automorphisms of $L$ fixing $K$ form a group, the \emph{Galois group} of $P$, and it permutes the roots of $P$; the theorem says that it is as large as the degree allows. Chapter~10 is built on this.

\begin{example}\label{ex:cuberoot}
The splitting field of $X^3-2$ over $\Q$. The roots are $\sqrt[3]2$, $\omega\sqrt[3]2$, $\omega^2\sqrt[3]2$ with $\omega=e^{2i\pi/3}$, so the splitting field in $\C$ is $L=\Q(\sqrt[3]2,\omega)$ (it contains $\omega=\omega\sqrt[3]2/\sqrt[3]2$, and conversely). $X^3-2$ is irreducible (Eisenstein at $2$), so $[\Q(\sqrt[3]2):\Q]=3$. The field $\Q(\sqrt[3]2)$ is real, so it does not contain $\omega$, a root of $X^2+X+1$: $[L:\Q(\sqrt[3]2)]=2$ and $[L:\Q]=6$. By Theorem~\ref{thm:splitunique}, $L$ has $6$ automorphisms. Each permutes the three roots and is determined by what it does to them, so every one of the $3!=6$ permutations of the roots comes from an automorphism: the Galois group is $\mathfrak S_3$.
\end{example}

A quadratic extension is built by adjoining a square root, exactly the operation that appears when two lines or circles intersect. The tower law therefore turns a geometric construction question into a restriction on algebraic degrees.

\section{Ruler and compass}

Identify the plane with $\C$. Start with the points $0$ and $1$. A \emph{ruler-and-compass step} draws a line through two points already constructed, or a circle whose centre is a constructed point and whose radius is the distance between two constructed points, and adds to the constructed points the intersections of the lines and circles drawn so far. A complex number is \emph{constructible} if it is obtained after finitely many steps. A real number is constructible if it is constructible as a point of $\C$; a regular $n$-gon is constructible if $e^{2i\pi/n}$ is (then so are all its vertices, by Proposition~\ref{prop:constrfield}).

\begin{proposition}\label{prop:constrfield}
The constructible numbers form a subfield $\mathcal C$ of $\C$, closed under complex conjugation and under square roots. A complex number $x+iy$ ($x,y$ real) is constructible iff $x$ and $y$ are.
\end{proposition}
\begin{proof}
We use four classical constructions, each with lines and circles through constructed points.

\emph{(i) Perpendicular bisector.} If $a\ne b$ are constructed, the circles centred at $a$ through $b$ and at $b$ through $a$ meet at two points, and the line through them is the perpendicular bisector of $[a,b]$; it meets $[a,b]$ at the midpoint.

\emph{(ii) Perpendiculars.} Let $\ell$ be a line through constructed points $a\ne b$, and $p$ a constructed point. The circle centred at $p$ with radius $|p-a|$, or $|p-b|$ if $p=a$, meets $\ell$ at $a$ (resp.\ $b$) and at its mirror image across the foot of the perpendicular from $p$. If these two points differ, (i) applied to them gives the perpendicular to $\ell$ through $p$. If they coincide, the line through $p$ and that point is already perpendicular to $\ell$. A parallel to $\ell$ through $p$ is a perpendicular to a perpendicular.

\emph{(iii) Coordinates.} The real axis is the line through $0$ and $1$, the imaginary axis is its perpendicular at $0$, and the circle centred at $0$ through $1$ meets it at $i$. The perpendiculars from $z=x+iy$ to the axes meet them at $x$ and $iy$; the circle centred at $0$ through $iy$ meets the real axis at $\pm y$. Conversely, from $x$ and $y$ one gets $iy$ (same circle), and $z$ is the intersection of the perpendiculars to the axes at $x$ and at $iy$. So $z$ is constructible iff $x$ and $y$ are; in particular $\bar z=x-iy$ is constructible with $z$.

\emph{(iv) Real field operations.} Let $a,b$ be constructible reals. The circle centred at $a$ with radius $|b|=|b-0|$ meets the real axis at $a\pm b$. For products, we may assume $a,b>0$ (signs are handled by $a\mapsto-a$, the second point of the circle centred at $0$ through $a$). The parallel to the line through $1$ and $ib$, drawn through $a$, meets the imaginary axis at $iab$, by the intercept theorem (Thales); then $ab$ by (iii). Likewise the parallel to the line through $a$ and $i$ drawn through $b$ meets the imaginary axis at $ib/a$.

So the constructible reals form a subfield $F$ of $\R$, and by (iii), $\mathcal C=\{x+iy:x,y\in F\}$. The formulas $(x+iy)(x'+iy')=(xx'-yy')+i(xy'+x'y)$ and $1/(x+iy)=(x-iy)/(x^2+y^2)$ show that $\mathcal C$ is a field.

\emph{Square roots.} If $a>0$ is constructible, the circle with diameter $[-1,a]$ (its centre is the midpoint, by (i)) meets the imaginary axis at $\pm i\sqrt a$: by the right angle in a semicircle, $\sqrt a^2=1\cdot a$. For $z\ne0$ in $\mathcal C$, $r=|z|$ is constructible (circle centred at $0$ through $z$), so are $\sqrt r$ and $u=z/r$, with $|u|=1$. If $u=-1$, $\sqrt z=i\sqrt r$. Otherwise $1+u\ne0$ points along the bisector of the angle between $1$ and $u$ (a rhombus), so $v=(1+u)/|1+u|$ satisfies $v^2=u$, and $\sqrt r\,v$ is a square root of $z$. All of these are obtained by field operations and real square roots, so they are in $\mathcal C$.
\end{proof}

\begin{theorem}[Wantzel]\label{thm:wantzel}
A complex number $z$ is constructible iff there is a chain of subfields of $\C$
\[
\Q=K_0\subset K_1\subset\dots\subset K_r,\qquad[K_j:K_{j-1}]=2,
\]
with $z\in K_r$. In particular, a constructible number is algebraic over $\Q$ and its degree is a power of $2$.
\end{theorem}
\begin{proof}
\emph{If.} By induction on $j$, $K_j\subseteq\mathcal C$. For $j=0$, $\mathcal C$ is a field containing $1$, so it contains $\Q$. Let $\alpha\in K_j\setminus K_{j-1}$; then $K_j=K_{j-1}(\alpha)$ by the tower law, and $\alpha$ is a root of some $X^2+bX+c$ with $b,c\in K_{j-1}$, so $\alpha=\frac{-b\pm\delta}2$ with $\delta^2=b^2-4c$. By Proposition~\ref{prop:constrfield}, $\alpha\in\mathcal C$, and $K_j=K_{j-1}(\alpha)\subseteq\mathcal C$ since $\mathcal C$ is a field.

\emph{Only if.} Let $z_0=0,z_1=1,z_2,\dots,z_n=z$ be the points of a construction, listed in the order they appear, with $z_k=x_k+iy_k$. Let $F_k=\Q(x_0,y_0,\dots,x_k,y_k)\subseteq\R$. Each line drawn passes through two points with coordinates in $F_{k-1}$, so has an equation $ux+vy=w$ with $u,v,w\in F_{k-1}$; each circle has an equation $x^2+y^2+ux+vy+w=0$ with $u,v,w\in F_{k-1}$ (its centre and squared radius are in $F_{k-1}$). The point $z_k$ is the intersection of two of these curves:
\begin{itemize}
\item two lines: Cramer's rule gives $x_k,y_k\in F_{k-1}$;
\item a line and a circle: if $v\ne0$, substitute $y=(w-ux)/v$ in the circle equation, to get a quadratic equation for $x_k$ over $F_{k-1}$, so $x_k\in F_{k-1}(\sqrt\Delta)$ for some $\Delta\in F_{k-1}$, $\Delta\ge0$, and then $y_k$ too (if $v=0$, exchange the roles of $x$ and $y$);
\item two circles: subtracting the equations gives a line (the centres differ, as two distinct circles with the same centre do not meet), and we are back to the previous case.
\end{itemize}
So $F_k=F_{k-1}(\sqrt\Delta)$ has degree $1$ or $2$ over $F_{k-1}$. Finally $z\in F_n(i)$, of degree $1$ or $2$ over $F_n$. Removing the steps of degree $1$ from $\Q=F_0\subseteq\dots\subseteq F_n\subseteq F_n(i)$ gives the chain. Then $[K_r:\Q]=2^r$ by the tower law, and the degree of $z$, which divides it (Corollary~\ref{cor:finitealg}), is a power of $2$.
\end{proof}

\paragraph{Why this works.}
Straightedge-and-compass constructions use intersections of lines and circles. Their coordinates come from linear or quadratic equations, never from a new irreducible cubic in one step. Starting with rational coordinates, each construction step therefore leaves the field unchanged or enlarges it by a square root. The tower law multiplies those degrees, so a constructible algebraic number must have degree a power of $2$. That is an \emph{obstruction}, not a complete test: an irreducible polynomial of degree $4$ may still have roots that cannot be constructed if the relevant field extension does not fit into a quadratic tower. The full theorem uses a quadratic \emph{tower}, not merely a degree bound.

\begin{figure}[ht]
\centering
\fig{ch08-pentagon}
\caption{Left: constructing the regular pentagon. As $2\cos72^\circ=\frac{\sqrt5-1}2$, the circle centred at $-\frac12$ through $i$ (radius $\frac{\sqrt5}2$) meets the real axis at $\frac{\sqrt5-1}2$; the perpendicular bisector of the segment from $0$ to that point is the line $x=\cos72^\circ$, which cuts the unit circle at two vertices. Right: the subfields of $\Q(\sqrt2,\sqrt3)$, with the degrees of the extensions.}
\label{fig:pentagon}
\end{figure}

\begin{corollary}[the classical problems]\label{cor:classical}
\begin{enumerate}[label=(\alph*)]
\item (Doubling the cube) $\sqrt[3]2$ is not constructible.
\item (Trisection) The angle $60^\circ$, which is constructible, cannot be trisected: $\cos20^\circ$ is not constructible.
\item (Polygons) The regular pentagon is constructible; the regular heptagon is not. More generally, if $p$ is a prime and the regular $p$-gon is constructible, then $p-1$ is a power of $2$.
\end{enumerate}
\end{corollary}
\begin{proof}
(a) $X^3-2$ is irreducible (Eisenstein at $2$), so $\sqrt[3]2$ has degree $3$, which is not a power of $2$.

(b) $e^{i\pi/3}=\frac12+i\frac{\sqrt3}2$ is constructible. By $\cos3\theta=4\cos^3\theta-3\cos\theta$ with $\theta=20^\circ$, $u=2\cos20^\circ$ satisfies
\[
u^3-3u-1=0.
\]
A rational root of this monic integer polynomial would be an integer dividing $1$ (Volume~1, Chapter~7), and $\pm1$ are not roots, so the cubic has no root in $\Q$; a reducible cubic has a factor of degree $1$, so it is irreducible, and $u$ has degree $3$. So $\cos20^\circ=u/2$ is not constructible, and neither is $e^{i\pi/9}$ (its real part would be, by Proposition~\ref{prop:constrfield}).

(c) Let $\zeta=e^{2i\pi/5}$. Dividing $\zeta^4+\zeta^3+\zeta^2+\zeta+1=0$ by $\zeta^2$, $t=\zeta+\zeta^{-1}=2\cos72^\circ$ satisfies $t^2+t-1=0$, so $t=\frac{\sqrt5-1}2$, which is constructible, and so is $\zeta=\frac t2+i\sqrt{1-t^2/4}$ (Figure~\ref{fig:pentagon}). For a prime $p$, $\zeta=e^{2i\pi/p}$ is a root of $\Phi_p=1+X+\dots+X^{p-1}$, irreducible over $\Q$ (Volume~1, Chapter~7, an exercise on Eisenstein's criterion); so $\zeta$ has degree $p-1$, which must be a power of $2$ if $\zeta$ is constructible. For $p=7$, $p-1=6$.
\end{proof}

\begin{remark}
Two classical results complete the picture; we quote them without proof here. \emph{Squaring the circle} asks to construct $\sqrt\pi$; that is impossible because $\pi$ is not even algebraic (Lindemann, 1882). And the converse of (c) holds: if $p=2^{2^k}+1$ is a (Fermat) prime, the $p$-gon is constructible. Gauss found the $17$-gon in 1796, aged 18; Chapter~10 proves the general statement with Galois theory, because ``degree a power of $2$'' alone is not enough to be constructible.
\end{remark}

\begin{beauty}
Wantzel turned a geometric question (``is there a construction?'') into an arithmetic one (``what is this degree?''), by looking not at the numbers but at the fields they generate. The same move proves that there is no formula in radicals for the roots of a general quintic (Chapter~10): there too, the obstruction is a structure attached to a field extension. The algebra of $\Q(\alpha)\cong\Q[X]/(m_\alpha)$ is also what computer algebra systems use to compute exactly with algebraic numbers, and the finite fields $\mathbb F_p[X]/(P)$ of the next chapter, built by the same Kronecker construction, are the arithmetic of error-correcting codes and of AES.
\end{beauty}

\begin{theorem}[primitive element]\label{thm:primitiveelement}
Let $K\subseteq\C$ be a field and $\alpha,\beta\in\C$ algebraic over $K$. There is $c\in\Q$ such that $K(\alpha,\beta)=K(\alpha+c\beta)$. Consequently every finite extension $L/K$ of subfields of $\C$ is of the form $K(\gamma)$.
\end{theorem}
\begin{proof}
Let $f=m_\alpha$ and $g=m_\beta$ (over $K$), with roots $\alpha=\alpha_1,\dots,\alpha_m$ and $\beta=\beta_1,\dots,\beta_n$ in $\C$, distinct by Lemma~\ref{lem:sep}. For $j\ge2$ and any $i$, the equation $\alpha_i+c\beta_j=\alpha+c\beta$ has at most one solution $c$. Since $\Q$ is infinite, choose $c\in\Q$ avoiding these finitely many values. If $n=1$, then $\beta\in K$ and $c=0$ works. If $n\ge2$, the value $c=0$ is excluded (take $i=1$), so $c\ne0$. Put $\gamma=\alpha+c\beta$ and $h(X)=f(\gamma-cX)\in K(\gamma)[X]$.

Then $g(\beta)=0$ and $h(\beta)=f(\alpha)=0$. If $g(\beta_j)=h(\beta_j)=0$, then $\gamma-c\beta_j=\alpha_i$ for some $i$, so $j=1$ by the choice of $c$. The only common complex root of $g$ and $h$ is $\beta$, and it is a simple root of $g$. Let $D$ be the monic gcd of $g$ and $h$ in $K(\gamma)[X]$. Then $D$ divides $g$ and $h$, so its complex roots are exactly the common roots, that is only $\beta$, and simply. Since $D$ splits over $\C$, $D=X-\beta$. The coefficients of $D$ lie in $K(\gamma)$, so $\beta\in K(\gamma)$, and then $\alpha=\gamma-c\beta\in K(\gamma)$. Hence $K(\alpha,\beta)=K(\gamma)$.

For a finite extension, write $L=K(\alpha_1,\dots,\alpha_n)$ and induct on $n$: $K(\alpha_1,\alpha_2)=K(\gamma_2)$, then $K(\gamma_2,\alpha_3)=K(\gamma_3)$, and so on.
\end{proof}

\section{Exercises}

\begin{exercise}[$\star$]\label{ex:sqrt23}
Let $\alpha=\sqrt2+\sqrt[3]2$. Show that $[\Q(\sqrt2,\sqrt[3]2):\Q]=6$, that $\sqrt2=\dfrac{\alpha^3+6\alpha-2}{3\alpha^2+2}$, and deduce that $\alpha$ has degree $6$ over $\Q$.
\end{exercise}

\begin{exercise}[$\star$]\label{ex:inverse}
In $\Q(\sqrt[3]2)$, write $\dfrac1{1+\sqrt[3]2}$ and $\dfrac1{1+\sqrt[3]2+\sqrt[3]4}$ in the form $a+b\sqrt[3]2+c\sqrt[3]4$ with $a,b,c\in\Q$.
\end{exercise}

\begin{exercise}[$\star\star$]\label{ex:classical}
Call an angle $\theta$ constructible if $e^{i\theta}$ is.
\begin{enumerate}[label=(\alph*)]
\item Show that if $\theta$ and $\varphi$ are constructible, so are $\theta\pm\varphi$ and $\theta/2$. Deduce that $3^\circ$ is constructible (use $72^\circ$ and $60^\circ$).
\item Show that the angle of $n$ degrees, $n\in\Z$, is constructible iff $3\mid n$. In particular the angle $72^\circ$ can be trisected, but the angle $1^\circ$ is not constructible.
\end{enumerate}
\end{exercise}

\begin{exercise}[$\star\star$]\label{ex:polygons}
\begin{enumerate}[label=(\alph*)]
\item Show that $u=2\cos\frac{2\pi}7$ satisfies $u^3+u^2-2u-1=0$, and that $2\cos\frac{2\pi}9$ satisfies $u^3-3u+1=0$. Deduce that the regular $7$-gon and $9$-gon are not constructible.
\item Show that $\cos\frac{2\pi}{17}$ has degree $8$ over $\Q$. (This is consistent with constructibility, but does not prove it: see Chapter~10.)
\end{enumerate}
\end{exercise}

\begin{exercise}[$\star\star$]\label{ex:algfield}
Show that the complex numbers algebraic over $\Q$ form a subfield $\overline{\Q}$ of $\C$, that $\overline{\Q}$ is countable, and that a complex root of a non-zero polynomial with coefficients in $\overline{\Q}$ lies in $\overline{\Q}$. Deduce that transcendental numbers exist.
\end{exercise}

\begin{exercise}[$\star\star$]\label{ex:splitdeg}
Find the degree over $\Q$ of the splitting field of $X^4-2$, and of $X^p-1$ for a prime $p$.
\end{exercise}

\begin{exercise}[$\star\star$]\label{ex:oddeg}
\begin{enumerate}[label=(\alph*)]
\item Show that if $[L:K]$ is prime, then $L=K(\alpha)$ for every $\alpha\in L\setminus K$.
\item Show that if $[K(\alpha):K]$ is odd, then $K(\alpha^2)=K(\alpha)$.
\end{enumerate}
\end{exercise}

\begin{exercise}[$\star\star\star$]\label{ex:primitive}
(Primitive element theorem.) Let $K\subseteq\C$ be a field and $\alpha,\beta\in\C$ algebraic over $K$. Show that there is $c\in\Q$ such that $K(\alpha,\beta)=K(\alpha+c\beta)$. Deduce that every finite extension $L/K$ of subfields of $\C$ is of the form $K(\gamma)$.
\end{exercise}

\section{Hints}

\begin{hint}{ex:sqrt23}
Both $2$ and $3$ divide the degree. For the formula, expand $(\alpha-\sqrt2)^3=2$ and collect the terms in $\sqrt2$.
\end{hint}
\begin{hint}{ex:inverse}
$(1+a)(1-a+a^2)=1+a^3$ and $(a-1)(1+a+a^2)=a^3-1$.
\end{hint}
\begin{hint}{ex:classical}
(a) $e^{i(\theta\pm\varphi)}=e^{i\theta}e^{\pm i\varphi}$, and square roots. (b) Bézout: if $3\nmid n$, $\gcd(n,360)$ is a combination of $n$ and $360$.
\end{hint}
\begin{hint}{ex:polygons}
(a) Divide $\Phi_7(\zeta)=0$ by $\zeta^3$ and write everything in terms of $\zeta+\zeta^{-1}$; for the $9$-gon, use $\cos3\theta$. (b) $\zeta=e^{2i\pi/17}$ is a root of $X^2-2\cos\frac{2\pi}{17}X+1$.
\end{hint}
\begin{hint}{ex:algfield}
$\alpha\pm\beta$, $\alpha\beta$, $\alpha^{-1}$ lie in $\Q(\alpha,\beta)$, a finite extension. For the last part, a tower $\Q\subseteq\Q(a_0,\dots,a_n)\subseteq\Q(a_0,\dots,a_n,z)$. Count polynomials with rational coefficients.
\end{hint}
\begin{hint}{ex:splitdeg}
The splitting field of $X^4-2$ is $\Q(\sqrt[4]2,i)$. For $X^p-1$, use $\Phi_p$.
\end{hint}
\begin{hint}{ex:oddeg}
(a) Corollary~\ref{cor:towerdiv}. (b) What is $[K(\alpha):K(\alpha^2)]$ at most?
\end{hint}
\begin{hint}{ex:primitive}
Choose $c$ so that the numbers $\alpha_i+c\beta_j$ (all roots of the two minimal polynomials) avoid $\alpha+c\beta$ except for $(i,j)=(1,1)$. Then compute the gcd of $m_\beta(X)$ and $m_\alpha(\gamma-cX)$ in $K(\gamma)[X]$, where $\gamma=\alpha+c\beta$.
\end{hint}

\section{Solutions}

\begin{solution}{ex:sqrt23}
\idea A degree is divisible by the degree of every intermediate field; to identify $\Q(\alpha)$ with the whole field, express the generators in terms of $\alpha$.
\writeup Let $L=\Q(\sqrt2,\sqrt[3]2)$. It contains $\Q(\sqrt2)$ and $\Q(\sqrt[3]2)$, of degrees $2$ and $3$, so $6\mid[L:\Q]$ (Corollary~\ref{cor:towerdiv}). Also $[L:\Q(\sqrt[3]2)]\le2$ ($\sqrt2$ is a root of $X^2-2$), so $[L:\Q]\le6$. Hence $[L:\Q]=6$.

From $\sqrt[3]2=\alpha-\sqrt2$: $2=(\alpha-\sqrt2)^3=\alpha^3-3\sqrt2\,\alpha^2+6\alpha-2\sqrt2$, so
\[
\sqrt2\,(3\alpha^2+2)=\alpha^3+6\alpha-2.
\]
As $3\alpha^2+2>0$, the formula follows. So $\sqrt2\in\Q(\alpha)$, then $\sqrt[3]2=\alpha-\sqrt2\in\Q(\alpha)$, and $\Q(\alpha)=L$ has degree $6$: $\alpha$ has degree $6$. (Its minimal polynomial, computed numerically, is $X^6-6X^4-4X^3+12X^2-24X-4$.)
\end{solution}

\begin{solution}{ex:inverse}
\idea Put $\alpha=\sqrt[3]{2}$, so that $\alpha^3=2$. An inverse in the degree-three field $\Q(\alpha)$ can be found by multiplying a denominator by a suitable quadratic expression. Here the two classical factorizations $1+x^3$ and $x^3-1$ do the work immediately.

\writeup The minimal polynomial $X^3-2$ is irreducible over $\Q$ (Eisenstein at $2$, or the rational-root test). Thus $1,\alpha,\alpha^2$ form a basis of $\Q(\alpha)$ over $\Q$. In particular, any answer has a unique expression
\[
a+b\alpha+c\alpha^2,\qquad a,b,c\in\Q.
\]

For the first denominator, use
\[
(1+\alpha)(1-\alpha+\alpha^2)=1+\alpha^3=3.
\]
Since $3\ne0$, division gives
\[
\boxed{\frac1{1+\sqrt[3]{2}}
=\frac13-\frac13\sqrt[3]{2}+\frac13\sqrt[3]{4}}.
\]
The identity can be checked by multiplying the right side by $1+\alpha$, after which all non-constant terms cancel and $1+\alpha^3=3$ remains.

For the second denominator, notice the exact factorisation
\[
(\alpha-1)(1+\alpha+\alpha^2)=\alpha^3-1=1.
\]
Therefore
\[
\boxed{\frac1{1+\sqrt[3]{2}+\sqrt[3]{4}}
=\sqrt[3]{2}-1}.
\]
This is also in the required form: $a=-1$, $b=1$, and $c=0$. Again the product with the given denominator is $1$, so no approximation of cube roots is involved.

Both computations illustrate the quotient-field viewpoint: an inverse is a polynomial in $\alpha$ whose product with the denominator is congruent to $1$ modulo the minimal polynomial $X^3-2$.
\end{solution}

\begin{solution}{ex:classical}
\idea Constructible angles form a group closed under halving; then $3^\circ$ is the gcd of what we can do and of what is forbidden.
\writeup (a) $e^{i(\theta\pm\varphi)}=e^{i\theta}\cdot(e^{i\varphi})^{\pm1}\in\mathcal C$, and $e^{i\theta/2}$ is a square root of $e^{i\theta}$, in $\mathcal C$ by Proposition~\ref{prop:constrfield} (it is $\pm$ the one we construct, and $-1\in\mathcal C$). $72^\circ$ (pentagon) and $60^\circ$ are constructible, so $12^\circ=72^\circ-60^\circ$, $6^\circ$ and $3^\circ$ are.

(b) If $3\mid n$, $n^\circ$ is a multiple of $3^\circ$, so it is constructible by (a). Conversely, suppose that $n^\circ$ is constructible with $3\nmid n$. Then $g=\gcd(n,360)$ is prime to $3$, and $g=un+360v$ with $u,v\in\Z$, so $g^\circ=u\cdot n^\circ+v\cdot360^\circ$ is constructible. Also $\gcd(g,3)=1=sg+3t$, so $1^\circ$ is constructible, hence $20^\circ=20\cdot1^\circ$, contradicting Corollary~\ref{cor:classical}(b). Thus $24^\circ$ is constructible ($3\mid24$) and $72^\circ$ can be trisected, while $1^\circ$ is not constructible.
\end{solution}

\begin{solution}{ex:polygons}
\idea A regular $n$-gon gives the complex number $\zeta_n=e^{2\pi i/n}$. The sum $\zeta_n+\zeta_n^{-1}=2\cos(2\pi/n)$ is real and often has a much smaller minimal polynomial. Find that polynomial first; only then use the degree obstruction to constructibility.

\writeup \emph{(a) The heptagon.} Let $\zeta=e^{2\pi i/7}$ and $u=\zeta+\zeta^{-1}$. Since
\[
1+\zeta+\zeta^2+\cdots+\zeta^6=0,
\]
division by $\zeta^3$ gives
\[
(\zeta^3+\zeta^{-3})+(\zeta^2+\zeta^{-2})
+(\zeta+\zeta^{-1})+1=0.
\]
We can express everything in $u$:
\[
\zeta^2+\zeta^{-2}=u^2-2,\qquad
\zeta^3+\zeta^{-3}=u^3-3u.
\]
Substituting and simplifying gives
\[
u^3+u^2-2u-1=0.
\]
The only possible rational roots of this monic integer cubic are $1$ and $-1$, and neither is a root. A reducible cubic over $\Q$ has a rational root, so this polynomial is irreducible. Hence $[\Q(u):\Q]=3$.

\emph{The nonagon.} Set $v=2\cos(2\pi/9)$ and $\theta=2\pi/9$. The triple-angle identity
\[
2\cos3\theta=(2\cos\theta)^3-3(2\cos\theta)
\]
gives $-1=v^3-3v$, since $3\theta=2\pi/3$. Thus
\[
v^3-3v+1=0.
\]
Again neither $1$ nor $-1$ is a root, so the cubic is irreducible and $[\Q(v):\Q]=3$.

The real part of a constructible complex point is constructible (Proposition~\ref{prop:constrfield}). If the regular heptagon or nonagon were constructible, its corresponding $u$ or $v$ would be constructible. Wantzel's theorem would then force its degree over $\Q$ to be a power of $2$. Degree $3$ is not, so neither polygon is constructible.

\emph{(b) The regular $17$-gon.} Let $\zeta=e^{2\pi i/17}$ and $c=\cos(2\pi/17)$. The cyclotomic polynomial $\Phi_{17}=1+X+\cdots+X^{16}$ is irreducible over $\Q$ by Eisenstein applied after substituting $X+1$ (as proved in Volume~1). Hence
\[
[\Q(\zeta):\Q]=16.
\]
Because $2c=\zeta+\zeta^{-1}$, the field $\Q(c)$ lies inside $\Q(\zeta)$. And $\zeta$ satisfies
\[
X^2-2cX+1=0
\]
over $\Q(c)$. It cannot belong to $\Q(c)$, which is a real field, because $\zeta$ is non-real. Therefore $[\Q(\zeta):\Q(c)]=2$. The tower law yields
\[
[\Q(c):\Q]=\frac{16}{2}=\boxed{8}.
\]
This degree passes the \emph{necessary} power-of-two test. By itself it is not a construction of the $17$-gon: that needs a suitable tower of quadratic extensions.
\end{solution}

\begin{solution}{ex:algfield}
\idea Replace an algebraic element by the finite field extension it generates. A finite collection of algebraic elements still generates a finite extension, which proves closure under field operations. Counting polynomials and their roots then gives countability.

\writeup Let
\[
\overline{\Q}=\{\alpha\in\C:\alpha\text{ is algebraic over }\Q\}.
\]
Clearly $0,1\in\overline{\Q}$. Suppose $\alpha,\beta\in\overline{\Q}$. Their simple extensions $\Q(\alpha)/\Q$ and $\Q(\beta)/\Q$ are finite. The tower law shows
\[
[\Q(\alpha,\beta):\Q]
=[\Q(\alpha,\beta):\Q(\alpha)]\,[\Q(\alpha):\Q]
<\infty,
\]
because $\beta$ satisfies its original polynomial over $\Q$, hence also over the larger coefficient field $\Q(\alpha)$.

The field $\Q(\alpha,\beta)$ contains $\alpha+\beta$, $\alpha-\beta$, $\alpha\beta$, and $\alpha/\beta$ when $\beta\ne0$. Every element of a finite extension is algebraic over the base field (Corollary~\ref{cor:finitealg}). These four expressions therefore lie in $\overline{\Q}$. This proves that $\overline{\Q}$ is a subfield of $\C$.

Now let $P(X)=a_dX^d+\cdots+a_0$ be a non-zero polynomial with coefficients in $\overline{\Q}$, and suppose $P(\gamma)=0$ for some $\gamma\in\C$. Put $M=\Q(a_0,\ldots,a_d)$. Repeated application of the preceding finite-extension argument makes $M/\Q$ finite. The equation $P(\gamma)=0$ says that $\gamma$ is algebraic over $M$, so $[M(\gamma):M]<\infty$. The tower law gives
\[
[M(\gamma):\Q]<\infty.
\]
Thus $\gamma$ is algebraic over $\Q$ and belongs to $\overline{\Q}$. In particular, every polynomial over $\overline{\Q}$ splits into linear factors \emph{within} $\overline{\Q}$, since its complex roots are all algebraic. So $\overline{\Q}$ is algebraically closed.

Finally, for each fixed degree $d$, there are countably many polynomials in $\Q[X]$ of degree $d$, since their coefficients form a subset of the countable set $\Q^{d+1}$. Taking the union over all degrees still gives countably many non-zero polynomials. Each has only finitely many complex roots. Therefore their union of root sets, which is exactly $\overline{\Q}$, is countable.

But $\C$ is uncountable (already $\R$ is uncountable, by Cantor's diagonal argument). Hence
\[
\C\setminus\overline{\Q}\ne\varnothing,
\]
indeed it is uncountable. Every element of this complement is transcendental over $\Q$. This counting proof establishes their existence without needing to exhibit a specific transcendental number.
\end{solution}

\begin{solution}{ex:splitdeg}
\idea A splitting field is obtained by adjoining \emph{all} roots, but often one root and one additional element already generate them. Identify those generators and use the tower law to multiply their extension degrees.

\writeup First consider
\[
f(X)=X^4-2.
\]
Let $\alpha=\sqrt[4]{2}$ be the positive real fourth root of $2$. The four roots of $f$ in $\C$ are
\[
\alpha,\quad-\alpha,\quad i\alpha,\quad-i\alpha.
\]
Adjoining $\alpha$ and $i$ therefore supplies all four roots, and conversely the quotient $(i\alpha)/\alpha=i$ shows that a field containing all roots must contain $i$. Thus the splitting field is exactly
\[
L=\Q(\alpha,i).
\]

The polynomial $X^4-2$ is irreducible over $\Q$ by Eisenstein's criterion at $2$: each non-leading coefficient is divisible by $2$, but the constant coefficient $-2$ is not divisible by $4$. Hence
\[
[\Q(\alpha):\Q]=4.
\]
The field $\Q(\alpha)$ is contained in $\R$, because $\alpha$ is real and its elements are rational functions of $\alpha$. So $i\notin\Q(\alpha)$. The polynomial $X^2+1$ has degree $2$ and has no root in this real field, giving
\[
[L:\Q(\alpha)]=2.
\]
By the tower law,
\[
\boxed{[L:\Q]=2\cdot4=8}.
\]
The Galois group therefore also has order $8$. Chapter~10 identifies it more precisely with the dihedral group of order $8$.

Now let $p$ be any prime and take $g(X)=X^p-1$. Its roots are
\[
1,\zeta,\zeta^2,\ldots,\zeta^{p-1},
\qquad \zeta=e^{2\pi i/p}.
\]
Thus the splitting field is $\Q(\zeta)$. The minimal polynomial of $\zeta$ is the $p$-th cyclotomic polynomial
\[
\Phi_p(X)=1+X+\cdots+X^{p-1}.
\]
For prime $p$ this polynomial is irreducible over $\Q$ (apply Eisenstein at $p$ to $\Phi_p(X+1)$). Its degree is $p-1$, so
\[
\boxed{[\Q(\zeta):\Q]=p-1}.
\]
For $p=2$ this reads $1$, as expected: $X^2-1$ already splits over $\Q$.
\end{solution}

\begin{solution}{ex:oddeg}
\idea The tower law leaves no room.
\writeup (a) $K\subsetneq K(\alpha)\subseteq L$, and the only intermediate fields are $K$ and $L$ (Corollary~\ref{cor:towerdiv}), so $K(\alpha)=L$.

(b) $\alpha$ is a root of $X^2-\alpha^2\in K(\alpha^2)[X]$, so $[K(\alpha):K(\alpha^2)]\in\{1,2\}$. It divides the odd number $[K(\alpha):K]$ (tower law), so it is $1$: $K(\alpha^2)=K(\alpha)$.
\end{solution}

\begin{solution}{ex:primitive}
\idea Recover $\beta$ from $\gamma=\alpha+c\beta$ as the only common root of two polynomials with coefficients in $K(\gamma)$; the gcd of these polynomials is then $X-\beta$, which has coefficients in $K(\gamma)$.
\writeup Let $f=m_\alpha$ and $g=m_\beta$ (over $K$), with roots $\alpha=\alpha_1,\dots,\alpha_m$ and $\beta=\beta_1,\dots,\beta_n$ in $\C$, distinct by Lemma~\ref{lem:sep} (a subfield of $\C$ contains $\Q$). For $j\ge2$ and any $i$, the equation $\alpha_i+c\beta_j=\alpha+c\beta$ has at most one solution $c$ (as $\beta_j\ne\beta$). Since $\Q$ is infinite, we can choose $c\in\Q$ avoiding these finitely many values. If $n=1$, $\beta\in K$ and $c=0$ works; if $n\ge2$, $c=0$ is one of the excluded values (take $i=1$), so $c\ne0$. Assume $n\ge2$, put $\gamma=\alpha+c\beta$ and $h(X)=f(\gamma-cX)\in K(\gamma)[X]$, a polynomial of degree $m$ since $c\ne0$.

Then $g(\beta)=0$ and $h(\beta)=f(\alpha)=0$. If $g(\beta_j)=h(\beta_j)=0$, then $\gamma-c\beta_j=\alpha_i$ for some $i$, so $j=1$ by the choice of $c$. So the only common complex root of $g$ and $h$ is $\beta$, and it is a simple root of $g$. Let $D$ be the monic gcd of $g$ and $h$ in $K(\gamma)[X]$. Then $D=ug+vh$ with $u,v\in K(\gamma)[X]$ and $D$ divides $g$ and $h$, so the complex roots of $D$ are exactly the common roots of $g$ and $h$, i.e.\ only $\beta$, and simply as $D\mid g$. Since $D$ splits over $\C$, $D=X-\beta$. Its coefficients are in $K(\gamma)$, so $\beta\in K(\gamma)$, then $\alpha=\gamma-c\beta\in K(\gamma)$. Hence $K(\alpha,\beta)\subseteq K(\gamma)\subseteq K(\alpha,\beta)$.

For the last part, write $L=K(\alpha_1,\dots,\alpha_n)$ (take a basis, say) and use induction on $n$: $K(\alpha_1,\alpha_2)=K(\gamma_2)$, then $K(\gamma_2,\alpha_3)=K(\gamma_3)$, and so on.
\end{solution}

\section*{Sources and further reading}

\begin{itemize}
\item P.~L.~Wantzel, \emph{Recherches sur les moyens de reconnaître si un problème de géométrie peut se résoudre avec la règle et le compas}, J.\ Math.\ Pures Appl.\ 2 (1837), 366--372. The impossibility of doubling the cube and trisecting the angle, in seven pages.
\item C.~F.~Gauss, \emph{Disquisitiones Arithmeticae}, 1801, Section~VII. Cyclotomy and the construction of the $17$-gon.
\item L.~Kronecker, \emph{Ein Fundamentalsatz der allgemeinen Arithmetik}, J.\ reine angew.\ Math.\ 100 (1887), 490--510. Adjoining a root by computing modulo a polynomial.
\item G.~Cantor, \emph{Über eine Eigenschaft des Inbegriffes aller reellen algebraischen Zahlen}, J.\ reine angew.\ Math.\ 77 (1874), 258--262. The algebraic numbers are countable.
\item F.~Lindemann, \emph{Über die Zahl $\pi$}, Math.\ Ann.\ 20 (1882), 213--225. $\pi$ is transcendental, so the circle cannot be squared.
\item E.~Artin, \emph{Galois Theory}, Notre Dame, 1942 (Dover reprint). Degrees and splitting fields done with linear algebra, the approach of this chapter.
\item I.~Stewart, \emph{Galois Theory}, 4th ed., CRC Press, 2015, Chapters~4--7. Ruler and compass, with many pictures and historical notes.
\end{itemize}
